\section{A finite-field formula for abelian Cayley towers}\label{sec:moments}
The preceding argument has a uniform form when the prime exceeds the
valency. Fix $m$ steps $s_i$ spanning $\mathbb F_p^d$ modulo $p$, with
$p>2m$. For a projective covector $L\in\mathbb P^{d-1}(\mathbb F_p)$
represented by $u$, define
\[
 r(L)=\min\{r\ge1:\textstyle\sum_{i=1}^m\langle u,s_i\rangle^{2r}
                       \ne0\ \text{in }\mathbb F_p\},
 \qquad A=2\sum_L r(L).
\]
Scaling $u$ preserves the first nonvanishing degree.

\begin{theorem}\label{thm:moments}
For every $n\ge1$,
\[
 \boxed{\displaystyle T_n=A\sum_{h=0}^{n-1}p^{(d-1)h}-dn.}
\]
In particular, the reductions of the steps modulo $p$ determine the
valuation at \emph{every} level. Each $r(L)$ is at most $m$.
\end{theorem}

\begin{proof}
Put $x_i=\langle u,s_i\rangle^2$ in $\mathbb F_p$. If their first $m$
power sums vanished, Newton's identities (with $1,\ldots,m$ invertible)
would give $\prod_i(Z-x_i)=Z^m$. This forces all projections zero and
contradicts spanning. Next take the truncated exponential coordinate
$1+T=\exp z$ modulo $z^{2m+1}$ over $\mathbb F_p$. All denominators used
are units, and the linear term is invertible. The eigenvalue becomes
\[
 -2\sum_{j=1}^m\frac{z^{2j}}{(2j)!}
                 \sum_i\langle u,s_i\rangle^{2j}
                 \pmod {z^{2m+1}} .
\]
Its first nonzero degree is $2r(L)$. The integral series has zero constant
and linear coefficients. Since $2r(L)\le2m\le p-1\le e$, the second
comparison rule proves that the orbit contributes exactly $2r(L)$.
There are $p^{(d-1)(h-1)}$ lifts of each projective residue direction to
conductor $p^h$: normalize the first unit coordinate to one and lift the
other $d-1$ coordinates freely. Formula~\eqref{eq:shells} now gives the claim.
\end{proof}

For example, for the square torus and $p\ge5$, the first moment is
$a^2+b^2$. For $p\equiv3\pmod4$ it never vanishes on a projective point.
For $p\equiv1\pmod4$ its two isotropic directions have nonzero fourth
moment $2a^4$. Consequently $A=2(p+1)$ and $A=2(p+3)$, respectively.
At $p=5$ this gives $T_n=4(5^n-1)-2n$. In rank one, Dion, Lei, Ray and
Vallières had already used the even moments and analyzed the distribution
of the one-variable Iwasawa invariant, including the two-loop case
\cite[Theorem 4.13 and Proposition 4.19]{DLRV}. The formula above counts
all residual projective directions and their lifts in arbitrary rank.

\paragraph{Sharpness in rank two.}
Suppose $m\ge2$ and $2m\mid p-1$. Take $g$ of order $2m$ modulo $p$
and lift the $m$ steps $(g^j,g^{-j})$, $0\le j<m$, to integers.
The first two span. In the direction $(a,b)$ the quadratic moment is
$2mab$, since the sums of $g^{2j}$ and $g^{-2j}$ vanish. Each coordinate
axis has all even moments through order $2m-2$ zero, and the $2m$-th moment is
nonzero. Thus $A=2(p-1)+4m$. This is the largest possible $A$ for
$m$ steps in rank two when their quadratic moment is not the zero polynomial:
a nonzero binary quadratic has at most two projective roots, each of
which can contribute at most $2m$. This proves sharpness for the stated
congruence class of primes. For example $p=13,m=6$ gives $A=48$.
